\documentclass[11pt]{article}
\usepackage{amsmath,amssymb,amsthm}
\newtheorem{theorem}{Theorem}
\title{Moment Cancellation by Binary Linear Sign Codes}
\author{}
\date{November 24, 2026}
\begin{document}
\maketitle

\paragraph{Status and scope.}
This is a complete elementary written lemma, not a Lean theorem and not
a lower bound for the original head model. It addresses a proposed
construction of many Boolean moment annihilators from a binary linear
code. Nonlinear sign codes, moving poles depending on the codeword,
and approximate rather than exact moment cancellation are outside its
scope.

Let $g\in\mathbb F_2^d$. Fix $H$ real poles $p_h$, nonzero real
coefficients $c_h$, and characters
$\chi_{\alpha_h}(g)=(-1)^{\alpha_h\cdot g}$ with
$\alpha_h\in\mathbb F_2^d$. Fixed signs in a code coset can be
absorbed into $c_h$. Put
\[
 F_g(t)=\sum_{h=1}^H\frac{c_h\chi_{\alpha_h}(g)}{t-p_h}.
\]
Poles may repeat. Suppose all codewords satisfy $k$ consecutive moment
cancellations:
\[
 \sum_h c_h\chi_{\alpha_h}(g)p_h^j=0
 \qquad(g\in\mathbb F_2^d, 0\leq j<k).
 \tag{1}
\]

\begin{theorem}
Let $r$ be the dimension over $\mathbb F_2$ of the span of the characters
whose aggregated rational coefficient functions are not identically
zero. Then
\[
 H\geq r(k+1),
 \qquad |\{F_g:g\in\mathbb F_2^d\}|\leq2^r.
 \tag{2}
\]
In particular, if the $2^d$ codewords give distinct strict query-sign
patterns on any fixed query set, then $H\geq d(k+1)$.
\end{theorem}
\begin{proof}
For each character $\alpha$, aggregate its terms as
\[
 f_\alpha(t)=\sum_{h:\alpha_h=\alpha}\frac{c_h}{t-p_h}.
\]
Multiply (1) by $\chi_\alpha(g)$ and average over the finite group.
Character orthogonality gives
\[
 \sum_{h:\alpha_h=\alpha}c_hp_h^j=0\qquad(0\leq j<k).
 \tag{3}
\]
Combine repeated poles within this character class and discard zero
combined coefficients. If $f_\alpha$ is nonzero and has $m$ remaining
distinct poles, then $m\geq k+1$. Otherwise the first $m$ equations in
(3) are a nonsingular Vandermonde system and force every coefficient
to be zero. Thus every active character costs at least $k+1$ original
terms. At least $r$ active characters are needed to span dimension $r$,
proving the first assertion.

The value of every active character is determined by its values on a
basis of their $r$-dimensional span. The identity
$F_g=\sum_\alpha\chi_\alpha(g)f_\alpha$ therefore gives at most
$2^r$ rational functions. Distinct strict sign words require distinct
functions, so $2^d$ such words force $r=d$.
\end{proof}

\paragraph{Consequence for a particular higher-order construction.}
The saved Catalan deformation exposes a $q$-dimensional leading query
polynomial after annihilating $k=q-1$ consecutive weighted monomial
moments. A $q$-bit binary linear sign code supplying all those
annihilators would require at least $q^2$ poles by this lemma.
Consequently that specific choice of code cannot supply a linear-pole
universal sign cube. This says nothing about a nonlinear subset of
Boolean annihilators. Even the elementary minimum-distance bound for
such a subset permits exponentially many words when $H$ is a sufficiently
large constant multiple of $q$.

\paragraph{Finite audit.}
The companion check verifies character projection and the sharp
construction consisting of $d$ independent finite-difference groups,
each with $k+1$ poles. Its pole count is exactly $d(k+1)$ and its
$2^d$ rational functions are distinct. The unrestricted claim is the
Fourier and Vandermonde proof above.
\end{document}
